%==============================================================================
\part{Execution}
%==============================================================================

\section{Budget allocation across heterogeneous sources}
\label{sec:allocation}

\subsection{Why allocation is forced}

The sources of \cref{obs:kind} differ in cost by orders of magnitude and in
kind of limit. A locally loaded artefact is effectively free and bounded by
memory; a SPARQL endpoint is bounded by a per-request timeout and an
unadvertised fair-use policy; a flat-file REST interface is bounded by an
explicit request rate. A plan that ignores this either underuses the cheap
sources or exceeds the expensive ones, and \cref{lst:script} does both. Since
the total request budget is finite and the sources' yields are different
functions of the effort spent on them, allocation is an optimisation problem
and not a scheduling convention.

\begin{definition}[Yield]
\label{def:yield}
The \emph{yield} of step $i$ under effort $e_i \geq 0$ is
$\gamma_i(e_i)$, where $\gamma_i : \Rset_{\geq 0} \to \Rset_{\geq 0}$ is
non-decreasing, continuously differentiable, strictly concave, and satisfies
$\gamma_i(0) = 0$. Effort is measured in requests.
\end{definition}

Strict concavity is the substantive assumption and it is the empirically
plausible one: the first requests to a source retrieve the common cases and
later requests retrieve progressively rarer ones. It fails for a step whose
result is all-or-nothing --- a single lookup that either succeeds or does not
--- and \cref{rem:concavity-fails} records the consequence.

\begin{definition}[Allocation problem]
\label{def:alloc}
Given a plan with $m$ steps, total budget $\Bud$, and per-step minimum costs
$c_i > 0$, the allocation problem is
\[
\max_{e \in \Rset^m_{\geq 0}} \sum_{i=1}^{m} \gamma_i(e_i)
\quad\text{subject to}\quad
\sum_{i=1}^{m} e_i \leq \Bud .
\]
\end{definition}

\subsection{A single shadow price}

\begin{theorem}[Water-filling allocation]
\label{thm:allocation}
The allocation problem of \cref{def:alloc} has a unique optimum $e^\ast$, and
there is a scalar $p^\ast \geq 0$ --- the shadow price of a request --- such
that for every $i$,
\[
e_i^\ast > 0 \implies \gamma_i'(e_i^\ast) = p^\ast,
\qquad
e_i^\ast = 0 \implies \gamma_i'(0) \leq p^\ast ,
\]
and $\sum_i e_i^\ast = \Bud$ whenever some $\gamma_i'(0) > 0$. Consequently
$e_i^\ast = (\gamma_i')^{-1}(p^\ast)$ on the support, and $p^\ast$ is the unique
root of $\sum_i \max\bigl(0, (\gamma_i')^{-1}(p)\bigr) = \Bud$.
\end{theorem}

\begin{proof}
The objective is a sum of strictly concave functions, hence strictly concave;
the feasible set is a compact convex simplex-like region. A strictly concave
function on a compact convex set attains a unique maximum, giving existence and
uniqueness.

Form the Lagrangian
$L(e,p,\lambda) = \sum_i \gamma_i(e_i) - p(\sum_i e_i - \Bud) + \sum_i \lambda_i e_i$
with $p \geq 0$ the multiplier of the budget constraint and $\lambda_i \geq 0$
those of the non-negativity constraints. Slater's condition holds --- the point
$e_i = \Bud/(2m)$ is strictly feasible --- so the KKT conditions are necessary
and sufficient. Stationarity gives
$\gamma_i'(e_i^\ast) - p^\ast + \lambda_i^\ast = 0$ for each $i$. Complementary
slackness gives $\lambda_i^\ast e_i^\ast = 0$. If $e_i^\ast > 0$ then
$\lambda_i^\ast = 0$ and $\gamma_i'(e_i^\ast) = p^\ast$. If $e_i^\ast = 0$ then
$\lambda_i^\ast \geq 0$ and $\gamma_i'(0) = p^\ast - \lambda_i^\ast \leq p^\ast$.

If some $\gamma_i'(0) > 0$ then the budget constraint is active at the optimum:
were $\sum_i e_i^\ast < \Bud$, adding $\delta > 0$ to that $e_i$ would remain
feasible and, by continuity of $\gamma_i'$ and $\gamma_i'(0)>0$, would strictly
increase the objective, contradicting optimality. Hence $\sum_i e_i^\ast = \Bud$.

Strict concavity makes each $\gamma_i'$ strictly decreasing and therefore
invertible on its range, so $e_i^\ast = (\gamma_i')^{-1}(p^\ast)$ where
positive, and $\sum_i \max(0,(\gamma_i')^{-1}(p))$ is continuous and strictly
decreasing in $p$ on the region where it is positive, hence has a unique root
at $\Bud$.
\end{proof}

\begin{corollary}[Nothing is dropped by a rule]
\label{cor:no-dropping-rule}
A step receives zero effort at the optimum if and only if
$\gamma_i'(0) \leq p^\ast$, that is, if and only if its marginal yield at the
first request is below the price. No step is excluded by a selection heuristic,
a priority ordering, or a source blacklist.
\end{corollary}

\begin{proof}
Immediate from the two implications of \cref{thm:allocation}, which are
exhaustive and mutually exclusive given strict concavity.
\end{proof}

\Cref{cor:no-dropping-rule} matters because the alternative --- a hand-written
priority list saying which sources to try first --- is what \cref{lst:script}
does implicitly through its statement order, and it is unfalsifiable: one
cannot ask of a priority list whether it is right. One can ask of $p^\ast$
whether the yields that determined it were estimated correctly, which is a
question with an answer.

\begin{proposition}[Sufficient total budget is not sufficient]
\label{prop:necessary-not-sufficient}
Let $c_i$ be the minimum cost at which step $i$ can return a non-empty payload.
Then $\Bud \geq \sum_i c_i$ is necessary for a plan to complete with all steps
$\vans$, and is not sufficient.
\end{proposition}

\begin{proof}
Necessity: each step must be issued at cost at least $c_i$ to return a
non-empty payload, and the costs are additive over steps and drawn from one
budget, so any completing execution spends at least $\sum_i c_i$.

Insufficiency: take $m=2$ with $c_1 = c_2 = 1$ and $\Bud = 2$. Let step $2$
consume the output of step $1$, and let $\cost_{\Src_2}$ be increasing in input
cardinality with $\cost_{\Src_2}(k) = k$. If step $1$ at cost $1$ returns a
payload of cardinality $3$, then step $2$ requires cost $3$, but only $1$
remains. The plan halts with step $2$ reporting $\vref$ by rule (R3) of
\cref{def:verdicts}, despite $\Bud \geq c_1 + c_2$. The obstruction is that
$c_2$ is a minimum over inputs and the realised input is not the minimising
one, and no static quantity bounds the realised input without executing step
$1$.
\end{proof}

\begin{remark}[The consequence for planning]
\label{rem:replanning}
\Cref{prop:necessary-not-sufficient} shows that a purely static allocation can
be infeasible at run time through no error of the allocator. Two responses are
available. The first is to re-solve \cref{thm:allocation} after each step with
the remaining budget and the realised cardinalities, which is correct and costs
one convex solve per step. The second is to allocate statically from upper
bounds on cardinality, which is cheap and wastes budget in proportion to how
loose the bounds are. The prototype takes the first, and \cref{sec:limits}
records that we have not established a regret bound for it against the
clairvoyant allocation --- the natural comparison, which we do not make.
\end{remark}

\begin{remark}[Where concavity fails]
\label{rem:concavity-fails}
A \textsf{lookup} step against a REST interface has an all-or-nothing yield:
one request retrieves the record and further requests retrieve nothing. Its
yield is a step function, neither strictly concave nor differentiable, and
\cref{thm:allocation} does not apply to it. The prototype handles such steps by
removing them from the optimisation, charging their fixed cost to the budget
first, and solving \cref{def:alloc} over the remainder --- which is correct
because their effort is not a decision variable, and which means the shadow
price governs only the steps that have a real allocation choice.
\end{remark}

\section{Structural interpolation}
\label{sec:interp}

\subsection{An instance of the defect}

The following is a documented observation and we restate it in full because the
argument of this section depends on the details rather than on the headline.

\begin{observation}[Two spellings, two answers]
\label{obs:2-397}
On 10 August 2026, two requests were issued to a public SPARQL endpoint of a
reaction knowledge base, with a JSON results media type. The two requests were
byte-identical except for one line. The first, using an abbreviated object
list, returned a count of $2$. The second, using two separate triple patterns,
returned a count of $397$. The ratio is $198.5$.
\end{observation}

\begin{lstlisting}[caption={Form A. Returns 2. The abbreviation is on line 6.},label=lst:comma,language=,keywordstyle=]
PREFIX rh:    <http://rdf.rhea-db.org/>
PREFIX rdfs:  <http://www.w3.org/2000/01/rdf-schema#>
PREFIX chebi: <http://purl.obolibrary.org/obo/CHEBI_>
SELECT (COUNT(DISTINCT ?r) AS ?n) WHERE {
  ?r rdfs:subClassOf rh:Reaction ; rh:status rh:Approved .
  ?r rh:side/rh:contains/rh:compound/rh:chebi ?a , ?o .
  VALUES ?aa { chebi:35238 chebi:37022 }
  VALUES ?ox { chebi:35179 chebi:36147 chebi:133294 }
  ?a rdfs:subClassOf* ?aa .
  ?o rdfs:subClassOf* ?ox .
}
\end{lstlisting}

\begin{lstlisting}[caption={Form B. Returns 397. Differs from \cref{lst:comma} only in that line 6 is written as two patterns.},label=lst:longhand,language=,keywordstyle=]
  ?r rh:side/rh:contains/rh:compound/rh:chebi ?a .
  ?r rh:side/rh:contains/rh:compound/rh:chebi ?o .
\end{lstlisting}

The two forms are semantically equivalent under the specification. An object
list is syntax: the grammar admits a comma-separated object list in a
predicate-object list, and the abbreviation is expanded during translation to
the abstract query --- before evaluation, and without a special case for
property paths \citep{Harris2013}. A conforming implementation must therefore
return the same count for both. The difference is not a specification
ambiguity.

\begin{remark}[The control is clean and the confound is named]
\label{rem:control}
Both forms were evaluated locally against two independent engines
\citep{RDFLib,Pyoxigraph} over a hand-checkable dataset, and both engines
returned the hand-computed answer for both spellings, so the equivalence is not
merely our reading of the grammar. Separately, the endpoint's loaded ontology
carried $204{,}585$ classes against $237{,}672$ in the corresponding published
download --- a difference of $33{,}087$, or $13.9\%$ --- so any comparison of
counts \emph{across} the two artefacts is confounded by snapshot skew. That
confound does not touch \cref{obs:2-397}, in which both counts come from the
same endpoint in the same session.
\end{remark}

\begin{remark}[We do not diagnose the service]
\label{rem:no-probing}
The cause lies in the endpoint's translation of the abbreviated form and we do
not attempt to determine it further. Establishing it would require probing a
third party's public service with requests designed to elicit internal
behaviour, which is not ours to do. Nothing in this section should be read as a
recommendation to do so, and the design consequence below does not depend on
knowing the cause.
\end{remark}

\subsection{The design consequence}

\Cref{obs:2-397} is an instance of a family. In each member, a plan author
writes a request that is correct under the target language's specification, a
deployment interprets one spelling differently from another, and the difference
is invisible because it is a difference between two strings the author regards
as the same string. Textual interpolation puts the author in exactly this
position, because the string that is sent is assembled at run time and never
inspected.

\begin{construction}[Longhand-only lowering]
\label{cons:longhand}
The lowering $\Low_\Src$ emits, for each abstract request, a canonical concrete
form fixed by the adapter, in which (i) every predicate-object list is expanded
to separate patterns, (ii) every bound result set enters through the target
language's own value-binding construct rather than by concatenation into the
pattern text, and (iii) no construct is emitted whose expansion interacts with
another construct in the sense of \cref{rem:independence-fails}.
\end{construction}

\begin{theorem}[Structural interpolation eliminates the defect family]
\label{thm:interpolation}
Let $\mathcal{D}$ be the family of defects in which two concrete requests with
the same denotation under the target specification are assigned different
denotations by a deployment, and whose difference lies in the spelling of a
construct that \cref{cons:longhand} canonicalises. Then no execution of a plan
under \cref{cons:longhand} exhibits a member of $\mathcal{D}$.
\end{theorem}

\begin{proof}
Fix a plan and an execution. Every concrete request issued is produced by
$\Low_\Src$, since by \cref{prin:leaves} no other source of concrete requests
exists in the language. By \cref{cons:longhand}, $\Low_\Src$ emits a single
canonical spelling for each abstract request: the expanded, longhand form with
bindings supplied structurally. Membership in $\mathcal{D}$ requires two
requests differing in a canonicalised spelling. But for any two abstract
requests with the same denotation, the emitted concrete forms are determined by
the canonical procedure and therefore agree wherever the canonicalisation
applies; and no non-canonical spelling is ever emitted, so no pair of emitted
requests differs in one. Hence no member of $\mathcal{D}$ is exhibited.
\end{proof}

\begin{remark}[Precisely what is and is not eliminated]
\label{rem:interp-scope}
\Cref{thm:interpolation} is a construction-by-exclusion result and its strength
is exactly the coverage of \cref{cons:longhand}. It eliminates the family
$\mathcal{D}$ --- differences attributable to a spelling the lowering
canonicalises. It does not eliminate: deployments that misinterpret the
canonical form itself; capability over-declaration
(\cref{rem:honesty-assumption}), which is a different defect with the same
symptom; or divergence between a service and a specification in constructs the
canonicalisation does not touch. The honest statement is that the language
removes the author's ability to make one class of mistake, by removing the
author's access to the string, and that this is a smaller claim than
correctness.
\end{remark}

\begin{proposition}[The elimination is not available to textual interpolation]
\label{prop:textual-cannot}
Let $I$ be an interface in which the caller supplies the concrete request as a
string. Then for any canonicalisation procedure $C$, there is a caller of $I$
whose issued request is not in the image of $C$.
\end{proposition}

\begin{proof}
The caller supplies an arbitrary string and $I$ transmits it. Since $C$'s image
is a proper subset of the concrete language --- proper because $C$
canonicalises, so at least two distinct strings map to one, and the
non-canonical one is outside the image --- a caller supplying that
non-canonical string issues a request outside the image. No property of $I$
prevents this, because $I$'s contract is to transmit what it is given.
\end{proof}

\Cref{prop:textual-cannot} is why \cref{prin:leaves} is a principle of the
language rather than a recommended practice: a recommendation is exactly the
thing \cref{prop:textual-cannot} says is unenforceable.

\section{Re-execution}
\label{sec:idempotence}

A plan is often run more than once --- to refresh a result, to recover from a
budget exhaustion, to check a suspicious answer. Whether re-execution is safe
is not obvious, because a plan mixes sources with different volatility and a
retention check compares against a threshold.

\begin{definition}[Stability]
\label{def:stable}
A source $\Src$ is \emph{stable over an interval} if the dataset it holds is
unchanged throughout. A translation map is stable if its relation is unchanged.
\end{definition}

\begin{theorem}[Characterisation of safe re-execution]
\label{thm:idempotent}
Let $\Plan$ be well-capability with total budget $\Bud$, executed twice over an
interval throughout which every source and every translation map it uses is
stable. Then the two executions assign the same verdict and the same payload to
every step if and only if both executions allocate every step an effort at
least its realised cost. If some step is allocated less in one execution than
its realised cost, the two executions may differ, and the difference is
confined to that step and its successors in $G(\Plan)$.
\end{theorem}

\begin{proof}
($\Leftarrow$) Suppose both executions allocate every step at least its
realised cost. We induct along the sequence order. For $\Step_1$,
$\beta_1 = ()$, so (R1) does not fire; $\Req(\rho_1) \subseteq \Capset(\Src_1)$ by
well-capability, so (R2) does not fire; the allocation covers the realised
cost, so neither (R3) nor (R4) fires; the request is issued, and by stability
the dataset is the same in both executions, so by faithfulness
(\cref{thm:lowering}) the extracted result equals
$\llbracket \rho_1 \rrbracket_D$ in both. Hence the verdict and payload agree.
For the inductive step, assume all predecessors of $\Step_i$ agree across the
two executions. Then the antecedent of (R1) has the same truth value in both,
and if it is false the same argument as the base case applies, using stability
of any translation map involved and the fact that a retention check compares a
fixed $r_\mu(S)$ against a fixed $\epsilon$, so it too agrees.

($\Rightarrow$) Suppose some step is allocated less than its realised cost in
some execution. Then in that execution rule (R3) or (R4) fires at that step,
yielding $\vref$ or $\vtime$; if the other execution allocates at least the
realised cost, that step yields $\vans$ or $\vempty$ there. The verdicts
differ, contradicting agreement. Hence agreement forces sufficient allocation
everywhere.

Confinement: by \cref{def:verdicts}, a step's verdict depends only on its own
source, capabilities, budget and the verdicts and payloads of its predecessors.
If all predecessors of $\Step_j$ agree and $\Step_j$ is not the differing step,
the same argument as above gives agreement at $\Step_j$. So the difference
propagates only along edges of $G(\Plan)$ out of the differing step.
\end{proof}

\begin{corollary}[Budget exhaustion is recoverable, data change is not]
\label{cor:rerun}
Re-running a plan that halted with $\vref$ or $\vtime$, under a larger budget
and over a stable interval, reproduces every verdict the first execution
assigned before the halt and may extend the plan further. Re-running across a
change in any source's dataset carries no such guarantee, and the language
records source snapshot identifiers in provenance so that the two cases are
distinguishable after the fact.
\end{corollary}

\begin{proof}
The first claim is \cref{thm:idempotent}($\Leftarrow$) applied to the prefix of
steps allocated sufficiently in both executions, which by hypothesis includes
every step completed before the halt. The second follows because stability is a
hypothesis of \cref{thm:idempotent} and the theorem's conclusion is not
asserted without it; the provenance recording is a design decision, not a
consequence.
\end{proof}
